% Related-work additions for Section II of Fleet-DT v2 (IEEE Access revision).
% Citation keys follow the [key] convention of sec2_patch.tex; refs_v2.bib holds the new keys.
% Each block names the subsection and the paragraph after which it goes.
% Reviewer items covered: (a) Sarr, Yang; (b) Rao, Djukic, Wei; (c) Leith, already in
% sec2_patch.tex; (d) Hoiland-Jorgensen; (e) Winstein; (f) 802.11ax/be; (g) Kaul, Yates,
% Li; (h) Kapteyn, NASEM; (i) Wu; (j) Tobagi.

%%% II-A: new paragraph after the paragraph on NDTs (after Almeida et al.)
Wu et al.\ survey digital twin networks, in which the twins of physical objects exchange data with these objects over a network [wu2021digital]. Kapteyn et al.\ formulate a DT as a probabilistic graphical model that couples the physical state, the observed data, the digital state, and the control actions over time, and update the digital state of an unmanned aerial vehicle in flight by Bayesian inference [kapteyn2021probabilistic]. The National Academies require continual verification, validation, and uncertainty quantification of a DT for the purpose of the DT [nasem2024foundational]. Fleet-DT states the fidelity of the DT of the network with the metrics of Definition~\ref{def:fidelity} and measures each metric against the ground truth of the emulator (Section~\ref{sec:method}).

%%% II-C: new paragraph after the paragraph on HCF/EDCA/HCCA, ETSI DCC, and available bandwidth
%%% (before the paragraph of sec2_patch.tex on the user priority and Leith et al.)
Tobagi and Kleinrock show that hidden terminals, stations that do not sense the transmissions of each other, degrade the throughput of carrier sense multiple access [tobagi1975packet]. Sarr et al.\ estimate the available bandwidth of an IEEE~802.11 link from the idle periods of the sender and of the receiver, the probability that these periods overlap, and the collision probability, and admit a flow when this estimate covers the rate of the flow [sarr2008bandwidth]. Yang and Kravets admit a flow when every node in the carrier-sensing range of its path has, by the fraction of idle channel time, the bandwidth that the flow consumes at this node [yang2005contention]. H{\o}iland-J{\o}rgensen et al.\ schedule the queues of the stations at the AP by airtime in the Linux kernel and obtain nearly equal airtime among the stations [hoiland2017ending]. Winstein et al.\ forecast the bytes that a cellular link delivers in the next intervals and send the bytes that a low quantile of this forecast allows, which bounds the risk of a delay beyond a target [winstein2013stochastic]. These works estimate or share the capacity per link, node, or station and serve flows with a rate. Fleet-DT predicts the capacity of the data window of the fleet in each step, admits the data of the step in the order of need of the GS, and bounds the long-run fraction of the steps with a missed admitted message (Section~\ref{sec:dt}); scenario S4 of Section~\ref{sec:method} emulates hidden vessels.

%%% II-D: new paragraph after the paragraph on IPACT and IEEE 802.16 (before network calculus)
Rao and Stoica divide the time into equal slots with loosely synchronized clocks above the IEEE~802.11 MAC and allocate the slots among the nodes with a distributed algorithm [rao2005overlay]. Djukic and Mohapatra synchronize commodity IEEE~802.11 nodes within microseconds and run TDMA schedules over several hops [djukic2009soft]. Wei et al.\ build RT-WiFi, a TDMA MAC on the timing synchronization function of IEEE~802.11, for wireless control loops [wei2013rtwifi]. IEEE~802.11ax lets the AP schedule uplink transmissions on resource units with trigger frames and the service periods of the stations with target wake time (TWT) [ieee80211ax2021, khorov2019tutorial], and IEEE~802.11be adds restricted TWT for latency-sensitive traffic [ieee80211be2024, adame2021tsn]. These mechanisms give slots or resource units to nodes, and each node selects the data of its slot. Fleet-DT admits the data of every vessel within one airtime budget per step in the order of need of the GS, on the EDCA of radios without trigger-based access (Section~\ref{sec:platform}), and Section~\ref{sec:method} compares Fleet-DT with an overlay TDMA and with the slotted variant of Definition~\ref{def:tdma}.

%%% II-E: append after the sentence of sec2_patch.tex that defines the AoI [kaul2012real]
Kaul et al.\ show for a first-come first-served queue that the average AoI grows when the update rate approaches zero or the service rate, so an intermediate update rate minimizes the average AoI [kaul2012real]. Yates et al.\ survey the AoI of queues, of scheduled sources, and of sources that share a channel [yates2021age]. Li et al.\ place the twins of mobile objects on edge servers, serve each query from the twin with the lowest AoI, and give a $1/4$ approximation for the static placement [li2024aoi]. These works minimize an average AoI over the updates of the sources. Fleet-DT admits the state of every vessel in every step and bounds the AoI of the state at the GS in every step that meets the hypothesis of Lemma~\ref{lem:delivery} (Corollary~\ref{cor:state}).
